\subsection{群芽的情形}

设 \((\mathbf{G}, e, \theta, m)\) 为李群芽，\(\Omega\) 为 m 的定义集。于是 \(\mathrm{T}(\Omega)\) 与 \(\mathrm{T}(\mathrm{G}) \times \mathrm{T}(\mathrm{G})\) 的一个开子集认同，而 \(\mathrm{T}(m)\) 是从 T(Ω) 到 T(G) 的解析映射。可以像第 2 节那样验证 (T(G), 0e, T(θ), T(m)) 是李群芽。G 与 T(G) 的乘法通常用乘法记号书写。T(G) 到 G 上的典范投影是李群芽态射。Te(m) 在 Te(G) 上的限制是 Te(G) 的向量空间加法。T(G) 的零截面是李群芽 G 到 T(G) 的一个李子群芽上的同构，我们把该李子群芽与 G 认同。若 f 是 G 到李群芽 H 的态射，则 T(f): T(G) → T(H) 是李群芽态射。

映射 \(\phi\colon(g,u)\mapsto gu\)（从 \(G\times T_{e}(G)\) 到 \(T(G)\)）是平凡向量丛 \(G\times T_{e}(G)\)（底空间为 \(G\)）到向量丛 \(T(G)\) 上的同构；因为 \(\phi\) 与 \(\phi^{-1}\) 都是解析的且都是向量丛态射，故只需引用 Differentiable and Analytic Manifolds, R, 7.2.1。（第 2 节的证明经过改动也可移用于此。）同构 \(\phi^{-1}\) 称为 \(T(G)\) 的左平凡化。映射 \((g,u)\mapsto ug\) 的逆同构称为右平凡化。

设 X 为 \(C^{r}\) 类流形，\(\psi\) 为 G 在 X 上的一个 \(C^{r}\) 类左作用律片段。于是 \(\mathrm{T}(\psi)\) 是一个 \(C^{r-1}\) 类的、\(\mathrm{T}(\mathrm{G})\) 在 T(X) 上的左作用律片段，且它是 \(\psi\) 的扩充。只要 gx 有定义，公式 (7)、(8) 与 (9) 便仍然有效。若 I 是 K 中包含 0 的开子集，若 \(\gamma: I \to G\) 是 \(C^{r}\) 类映射且满足 \(\gamma(0) = e\)，且若 \(a = \mathrm{T}_{0}(\gamma)1\)，则在 X 上定义的向量场 \(x \mapsto ax\) 就是 Differentiable and Analytic Manifolds, R, 8.4.5 意义下由映射 \((\lambda, x) \mapsto \gamma(\lambda)x\) 所定义的向量场。

